% !TEX TS-program = pdflatex
% Combined article for the three DuplexCascade pt-BR approaches.
% Supersedes distill/article/main.tex and mtp_like/article/mtp_duplex.tex as the
% single write-up; those remain as per-approach records.
\documentclass[conference]{IEEEtran}
\IEEEoverridecommandlockouts
\usepackage{cite}
\usepackage{amsmath,amssymb,amsfonts}
\usepackage{mathtools}
\usepackage{algorithmic}
\usepackage{graphicx}
\usepackage{textcomp}
\usepackage{booktabs}
\usepackage{multirow}
\usepackage[dvipsnames]{xcolor}
\usepackage[utf8]{inputenc}
\usepackage[T1]{fontenc}

\def\BibTeX{{\rm B\kern-.05em{\sc i\kern-.025em b}\kern-.08em
    T\kern-.1667em\lower.7ex\hbox{E}\kern-.125emX}}

\begin{document}

\title{Three ways to learn the DuplexCascade micro-turn protocol:\\
supervised fine-tuning, self-distillation, and a policy companion}

\author{
\IEEEauthorblockN{Felipe Penhorate Carvalho da Fonseca}
\IEEEauthorblockA{}
}

\maketitle

\begin{abstract}
Full-duplex spoken dialogue requires a conversational model to do two things at
once: produce answer content and control turn-taking---when to speak, when to
stop, when to acknowledge, when to yield. The DuplexCascade paradigm encodes both
in one language model by interleaving answer text with control tokens over
\emph{micro-turns}. Learning that protocol is cheap, but the fine-tune is
destructive: in our Brazilian-Portuguese implementation the plain-SFT model that
learns the tags with 98\% accuracy simultaneously loses most of its function
calling (100\%$\rightarrow$55.6\%) and general instruction following
(100\%$\rightarrow$41.7\%). This paper studies three ways to obtain the protocol
without that cost, on one shared pt-BR data pipeline, and evaluates them with one
shared harness. (i)~\textbf{SFT} is the original recipe, kept as the baseline and
as the fastest deployable model. (ii)~\textbf{Self-distillation} splits the
objective by token position: weighted cross-entropy on the tags exactly as
before, plus a per-token forward-KL to the frozen base model on every content
position. It needs neither reinforcement learning nor on-policy rollouts, and
with an offline teacher cache it costs no more than SFT; it matches the protocol
(96.4\% vs.\ 98\% tag accuracy) while holding function calling, general
instruction-following and knowledge benchmarks at base level. (iii)~An
\textbf{MTP-style protocol companion} separates the skills instead of balancing
them: a 0.8B sidecar predicts \emph{only} the control stream while an untouched
big model writes the answer, reaching 95.6\% tag-stream accuracy with 0.97/0.99
F1 on the two critical controller decisions (barge-in interruption and
backchannel) and supporting instruction-conditioned interrupts. All three are
then evaluated in \emph{free generation} with a Full-Duplex-Bench harness, and
against the upstream English DuplexCascade release re-measured on the same
harness, which reproduces its own published turn-taking numbers (ATA 0.896 vs.\
0.858 published). All three pt-BR variants beat that published accuracy
(ATA 0.955--0.973). The comparison also exposes what teacher-forced accuracy
cannot see: the distill model, despite the best capability retention, has the
weakest interruption relevance (1.59 vs.\ 3.56 and 4.36) and commits to the floor
$13\times$ more slowly (3.29\,s vs.\ 0.25\,s), because the KL anchors the
teacher's \emph{generation style} along with its capabilities. We further show
that the evaluated input stream must match the training stream---on the same model,
one ablation arm more than triples pause-handling TOR and another destroys
interruption recovery, and the two axes are separable---and that one FDB dimension
is unanswerable on this data by construction, since the system-backchannel tag
occurs zero times in training.
\end{abstract}

\begin{IEEEkeywords}
full-duplex dialogue, turn-taking, self-distillation, catastrophic forgetting,
parameter-efficient fine-tuning, MTP, VAD-free, evaluation
\end{IEEEkeywords}

\section{Introduction}
\IEEEPARstart{S}{poken} dialogue systems are moving from turn-based
(``push-to-talk'') interaction to \emph{full-duplex} conversation, where both
parties may speak at any time. The DuplexCascade
paper~\cite{duplexcascade} established that a language model can learn the
full-duplex protocol directly: its training data interleaves user and assistant
\emph{micro-turns} and supervises a set of control tokens (``user is speaking'',
``user finish speaking'', ``user interruption'', ``user backchannel'', ``user is
thinking'', ``system backchannel'') so that the same model that writes the answer
also emits the control stream. The model must therefore be fine-tuned---deeply---
for the turn-taking skill.

Learning the protocol is easy: the supervision is cheap, synthetic and
deterministic. The fine-tune is not. In our pt-BR implementation, the plain-SFT
model that learns the tags with 98\% accuracy simultaneously loses most of its
ability to call functions (100\%$\rightarrow$55.6\%) and to follow general
instructions (100\%$\rightarrow$41.7\%). This is catastrophic forgetting in its
most practical form: the model becomes a good turn-taker and a bad assistant.

This paper studies three answers to that problem, all built on the same pt-BR
duplex pipeline and all evaluated with the same harness:

\begin{itemize}
  \item \textbf{SFT} (Section~\ref{sec:sft})---the original recipe: QLoRA on the
        duplex micro-turn data with weighted, position-masked cross-entropy. It
        remains the best raw turn-taker and the fastest to deploy, and it is the
        baseline the other two are measured against.
  \item \textbf{Self-distillation} (Section~\ref{sec:distill})---the observation
        that the new skill and the capabilities to preserve are \emph{separable
        by token position}. The tags live at supervised positions; the content
        lives at the rest. Splitting the objective accordingly---weighted CE on
        the tags, per-token forward-KL to the frozen base everywhere else---
        retains the capabilities without RL and without rollouts.
  \item \textbf{An MTP-style protocol companion}
        (Section~\ref{sec:companion})---the inverse decomposition: a 0.8B sidecar
        learns \emph{only} the protocol and drives an untouched big model, which
        is the same division of labour as a Multi-Token-Prediction head except
        that the predicted vocabulary is the control protocol rather than the
        next answer token.
\end{itemize}

\emph{Is the protocol actually used?} The protocol metric these methods are
usually compared on is teacher-forced: at each supervised position the model is
handed the correct prefix and only its next-token choice is scored. That measures
whether the protocol was \emph{learned}; it says nothing about whether the model
\emph{decides} correctly when generating on its own, which is what a deployed
duplex system needs. We therefore evaluate all three, plus the untrained base as
a control, with a free-generation harness modelled on Full-Duplex-Bench
(FDB)~\cite{lin2025fdb}, and---to anchor the comparison---we re-measure the
upstream English \textsc{DuplexCascade} release with the same harness instead of
quoting it. That reference lands on its own published numbers, which is the
strongest available check that the harness measures what the benchmark measures.

Four findings follow from pairing the two evaluations, and they are the
contribution of Sections~\ref{sec:eval} and~\ref{sec:results}:

\begin{enumerate}
  \item \textbf{All three approaches reach the protocol; they differ in what it
        costs.} SFT and the companion reach the best turn-taking accuracy
        (ATA 0.971/0.973), distill is close (0.955) and all three exceed the
        published 0.858---but the companion also has the best interruption
        relevance (4.36) and the fastest commitment among the decomposed designs,
        while SFT is an order of magnitude faster to stop and respond (0.035\,s
        vs.\ 0.066\,s and 0.351\,s).
  \item \textbf{Teacher forcing hides a generation-style cost.} The distill model
        has the best capability retention and the worst free-generation reply
        quality (interruption relevance 1.59) with a $13\times$ slower floor
        commitment. The content-KL anchors the teacher's generation style
        together with its capabilities, and the teacher occasionally
        autocompleted the user's turn---an artifact teacher-forced accuracy
        (96.4\%) cannot detect.
  \item \textbf{The input stream is an experimental variable.} The models are
        trained on 1--7-token ASR pieces and $\sim$10-word turns. Two single-factor
        ablations show the axes are separable: multi-sentence turns more than
        triple pause-handling TOR (0.080$\rightarrow$0.307) while whole-sentence
        pieces destroy interruption recovery (0.995$\rightarrow$0.000). Any
        free-generation protocol benchmark must therefore report the input
        distribution it used.
  \item \textbf{A metric can be unanswerable by construction.} No model here can
        score the FDB backchannel dimension, because the tag it requires
        (\texttt{<|system backchannel|>}) occurs zero times in the training data.
        The upstream release, trained elsewhere, scores it---so the column
        measures a data gap, not a capability gap.
\end{enumerate}

\section{Background and Related Work}
\label{sec:background}

\subsection{DuplexCascade and the micro-turn protocol}
DuplexCascade~\cite{duplexcascade} keeps the intelligence of a text LLM by not
modelling audio: a streaming ASR feeds text to the LLM, which emits interleaved
answer text and control tokens, and a streaming TTS speaks the text. Turn-taking
is \emph{VAD-free} in the sense that no separate voice-activity detector decides
who speaks: the LLM itself emits control tokens based on the incoming text
stream. The dialogue is decomposed into \emph{micro-turns}---short interleaved
user and assistant fragments---so that the system can react before an utterance
is complete. We adopt this setting exactly, and our contributions are about
\emph{how the protocol is learned} and \emph{how it is measured}, not about the
architecture.

\subsection{Catastrophic forgetting and off-policy fine-tuning}
Sequential SFT is inherently off-policy: the model learns to imitate expert
actions on a fixed offline distribution, which prior work shows causes severe
forgetting when adapting to a new task or
domain~\cite{kirkpatrick2017,lifhoiem2017,mecklenburg2024}. The standard remedy
is to make the signal on-policy with reinforcement learning and an explicit
reward~\cite{chu2025,airef}. Rewards are often unavailable for a protocol skill,
and on-policy RL (PPO/GRPO) is expensive and brittle. This motivates the
distillation-based approaches we compare against plain SFT.

\subsection{Self-distillation}
\emph{Self-Distillation Fine-Tuning} (SDFT)~\cite{selfdistill2026} obtains
on-policy benefits from demonstrations alone: the model acts as a \emph{teacher}
conditioned on a demonstration and as a \emph{student} conditioned only on the
query, and training minimizes a KL between the two on the student's own
trajectories. It reduces forgetting substantially relative to SFT, but still
needs on-policy rollouts (roughly $4\times$ SFT wall-clock). Related
on-policy-distillation work reaches similar conclusions from different
directions~\cite{tml2025,airef,rlsrazor2025,bai2022,snell2022}. Our
self-distillation variant (Section~\ref{sec:distill}) is the degenerate case of
this family in which the demonstration is the training sequence itself, so the
teacher is simply the frozen base and no rollouts are required.

\subsection{Parameter-efficient fine-tuning}
We use LoRA~\cite{lora2022} and its quantized
variant~\cite{qlora2023} throughout, both for the 4B duplex models and for the
0.8B companion, with the toolchain of~\cite{unsloth}. PEFT matters here for a
second reason: because the student is \emph{base + adapter}, the teacher forward
pass is obtained by temporarily disabling the adapters, so the distillation
teacher costs no extra VRAM and no second checkpoint.

\subsection{Multi-token prediction and VAD-free turn-taking}
Multi-Token Prediction heads~\cite{deepseek} predict several future tokens from a
shared trunk with a lightweight module. Our companion
(Section~\ref{sec:companion}) borrows the \emph{architectural} idea---a small
module running alongside a large frozen model---but predicts a different
vocabulary: the duplex protocol. It is therefore closer in spirit to a
speculative or auxiliary head than to a pretraining objective. The
companion is trained on the same data and with the same distill objective as
Section~\ref{sec:distill}, which makes the two directly comparable.

\section{Shared infrastructure: the pt-BR micro-turn pipeline}
\label{sec:shared}

All three approaches consume one dataset and target one protocol vocabulary, so
we describe them once.

\subsection{Protocol vocabulary and supervision}
A dialogue is a strict alternation of user and assistant \emph{items}, each
rendered in ChatML and closed with \texttt{<|im\_end|>}; assistant items may carry
one control tag at their start. The tags are \texttt{<|user is speaking|>},
\texttt{<|user finish speaking|>}, \texttt{<|user is thinking|>},
\texttt{<|user interruption|>}, \texttt{<|user backchannel|>},
\texttt{<|system backchannel|>}, plus the input-only marker
\texttt{<|no voice|>} used for silence. Following the paper, loss is computed only
on assistant micro-turns and each tag carries a task-specific weight:
\texttt{user finish speaking} $\times 10$, \texttt{user interruption} $\times 5$,
\texttt{system backchannel} $\times 3$, \texttt{user backchannel} $\times 2$,
\texttt{user is speaking} and \texttt{user is thinking} $\times 1$.

\subsection{Data: source dialogues and micro-turn construction}
Two generations of source dialogues exist, and the difference matters for
evaluation (Section~\ref{sec:fdb_setup}). The first generator instructed
\emph{``respostas curtas do Assistente. Nada de textos longos''}, and the chunker
fixed every assistant micro-turn at exactly 10 tokens; the resulting models spoke
only in $\sim$10-token bursts. The continuation generator
(\texttt{build\_long\_dialogues.py}) mixes \emph{short} (weight 6), \emph{long}
(2) and \emph{mixed} (2) styles, so a fraction of assistant turns are 2--4-sentence
replies up to 220 words, and system chunking samples 10--48 tokens instead of a
constant. Both models in this paper train on the continuation data; the
companion trains on the same dialogues.

The deterministic micro-turn builder then converts each dialogue into items with
six phenomena: randomized micro-turn length (user chunks 1--7 tokens), natural
pauses (1--5 \texttt{<|no voice|>} items), user interruptions ($p=0.30$, with the
first interrupting chunk labelled \texttt{user interruption} and later ones
\texttt{user is speaking}), user backchannels ($p=0.01$, from a pt-BR
interjection list), system backchannels (opt-in, $p=0.20$ at sentence-final user
chunks) and user thinking pairs (1--20 \texttt{<|user is thinking|>} items closing
each completed system turn). Strict user/assistant alternation is asserted. The
generator \emph{oversamples interruptions} in FDB's sense in one respect worth
noting: it never samples the position where an interruption begins, so we treat
its 0.30 rate as an augmentation strength rather than a label prior.

\subsection{Measured turn sizes}
Because FDB's turn-taking gates are \emph{durations}, both sides' length
distributions are load-bearing. Measured over the training data (3{,}000
sequences):

\begin{itemize}
  \item \textbf{User pieces}: mean \textbf{2.9 words} per piece (median 3, p90 5),
        arriving as 3.7 pieces per turn---so a user turn is $\sim$10.5 words and
        contains $\sim$1.5 sentence-final pieces.
  \item \textbf{Assistant content bursts}: mean 6.1 words but \emph{median 5,
        p90 12, max 37}, with 50\% of bursts at ${\le}5$ words and 1\% above 20.
        The heavy right tail is the intended effect of the variable chunking and
        the longer source dialogues.
\end{itemize}

The consequence for evaluation is a mixture rather than a mean: a model trained
on longer bursts crosses FDB's ``formal turn'' threshold ($\ge$1\,s) with fewer
words, so the assistant-side length distribution moves the metric being reported.
We return to this in Section~\ref{sec:fdb_setup} and report both sides' lengths
with the results.

\subsection{Runtime}
The live stack is four components: \texttt{llama.cpp} serving the fine-tuned GGUF,
\texttt{faster-whisper} for streaming ASR, \texttt{pocket-tts} for streaming
speech synthesis, and a bridge that wires them and exposes a browser UI. Two
serving details are mandatory and were sources of silent failure: the server must
be started with \texttt{-sp} so the special tokens are returned by the completions
endpoint at all, and the bridge must post to the raw completions endpoint rather
than the chat endpoint, because the chat template mangles \texttt{<|...|>} tokens
appearing in assistant history.

\section{Approach A: supervised fine-tuning (the original recipe)}
\label{sec:sft}

\subsection{Method}
QLoRA with $r{=}16$, $\alpha{=}16$, dropout 0, applied to all linear modules,
plus \texttt{modules\_to\_save=["embed\_tokens","lm\_head"]} so that the seven new
special-token embeddings are trained. The base is
\texttt{Qwen/Qwen3-4B-Instruct-2507}. The objective is a weighted,
position-masked cross-entropy: labels are shifted by one, every user position and
every ChatML header is masked to $-100$, and each supervised tag token carries the
weight of Section~\ref{sec:shared}. Training runs in two stages: an initial
fine-tune, then a continuation fine-tune started from the \emph{merged} model on
a short+long response mixture. Starting the second stage from the merged model
rather than the base is required---from the base, the learned special-token
embeddings and the turn-taking pattern are re-randomized and lost.

\subsection{What it costs}
SFT is the fastest model to serve and, as Section~\ref{sec:results} shows, the
best raw turn-taker of the three. Its cost is the capability regression of
Section~\ref{sec:results}: it is the baseline that motivates the other two
approaches, not a straw man.

\section{Approach B: self-distillation by position}
\label{sec:distill}

\subsection{The observation}
The ``new skill'' (the turn-taking protocol) and the ``everything else'' (the
capabilities we want to preserve) are \emph{separable by token position}: the tags
live at supervised positions, the content lives at the other positions. We
therefore split the objective in two:

\begin{itemize}
  \item \textbf{Tags}: the weighted cross-entropy to the supervised tag tokens,
        exactly as in the plain SFT recipe.
  \item \textbf{Everything else}: a per-token forward-KL to the \emph{frozen base
        model} on content positions, so the student's distribution is anchored to
        the pretrained policy wherever there is no explicit tag supervision.
\end{itemize}

The teacher is the frozen original model itself conditioned on the same duplex
sequence, so no second model and no rollouts are needed. Because the student is
\emph{base + LoRA}, the teacher forward is obtained with zero extra VRAM by
temporarily disabling the adapters, and the teacher's top-$k$ next-token logits
are precomputed once (offline cache), which removes the second forward pass
entirely and makes the run no more expensive than SFT.

\subsection{Loss}
With $\hat{p}_i$ the student distribution at position $i$, $p^{\text{base}}_i$ the
frozen base distribution, $y_i$ the supervised tag and $w_i$ its weight,
\[
\begin{aligned}
\mathcal{L} ={}&
\underbrace{\frac{1}{N_{\text{sup}}}\sum_{i\in\text{tags}}
w_i\,\text{CE}(\hat{p}_i, y_i)}_{\text{tag cross-entropy}}\\
&+\;\lambda\,\underbrace{\frac{1}{N_{\text{sup}}}\sum_{i\in\text{content}}
D_{\text{KL}}\!\left(p^{\text{base}}_i \,\|\, \hat{p}_i\right)}_{\text{content
anchoring}},
\end{aligned}
\]
with $\lambda=1.0$ chosen by a short sweep over $\lambda\in\{0.1,0.3,1.0\}$ scored
on tag accuracy and KL drift.

\subsection{Practical details}
Two implementation details mattered. First, \emph{artifact masking}: in every
\texttt{user finish speaking} item the first two \emph{content} tokens of the
response are left unsupervised, so the model is not forced to reproduce the
teacher's opening tokens verbatim while the tag itself---predicted by the
position before it---stays fully supervised. Second, \emph{capability anchors}:
300 anchor prompts (150 function calling, 150 general QA) are mixed into the
training set with \texttt{is\_anchor} set, so the formats we intend to protect are
also anchored by the KL rather than only by the base weights.

\section{Approach C: an MTP-style protocol companion}
\label{sec:companion}

\subsection{Separation of concerns}
Instead of teaching the answer model the protocol, we train a \emph{small
companion} that predicts \emph{only} the protocol stream on behalf of the bigger
model, while the bigger model continues to do what it already does: produce answer
text. This is the same division of labour as an MTP
head~\cite{deepseek}---a lightweight module alongside the main model---except
that our companion predicts a different vocabulary: the duplex protocol, not the
next answer token. Where no control token applies it emits nothing (an empty
string that never enters the context). The companion is therefore the
\emph{controller} of the system: it decides turn starts and ends, barge-ins,
backchannels and silences, without any voice-activity detection.

\subsection{Vocabulary and supervision}
Let $\mathcal{T}$ be the companion's output vocabulary, a set of nine names:
\texttt{no\_tag}, \texttt{user\_is\_speaking}, \texttt{user\_finish\_speaking},
\texttt{user\_interruption}, \texttt{user\_backchannel},
\texttt{user\_is\_thinking}, \texttt{system\_backchannel},
\texttt{system\_take\_floor} and \texttt{system\_handover}, where
\texttt{no\_tag} is a label-only symbol that never enters the context. For every item $i$ of a dialogue the target is the tag
carried by the \emph{next} assistant item:
\[
\text{target}(i) = \begin{cases}
  \text{special}(\text{item}_{i+1}) & \text{item}_i \text{ is user}\\
  \texttt{no\_tag} & \text{item}_i \text{ is assistant}
\end{cases}
\]
Each dialogue becomes one sequence with labels masked everywhere except at item
boundaries. This gives the companion exactly the supervision the paper's model
received, but as a \emph{classification} of the protocol state rather than a
generation task.

\subsection{Controller augmentations}
The vanilla targets cover the paper's phenomena, but two behaviours are critical
for a controller and under-represented, so the prep adds $\sim$946k records on top
of the 10k main streams:

\paragraph{Streaming-aware supervision.} The runtime queries the companion while
user words are still arriving, so every multi-token user chunk also yields
partial-prefix records labelled \texttt{user\_is\_speaking}---except partials of
an interrupting chunk, labelled \texttt{user\_interruption}.

\paragraph{Barge-in records.} At random points inside every system turn the
current assistant item is rendered \emph{open} (speech in progress) followed by
(a)~the user's next real words $\rightarrow$ \texttt{user\_interruption},
(b)~\texttt{<|no voice|>} $\rightarrow$ \texttt{no\_tag} (silence during speech
means keep talking), or (c)~a sampled pt-BR backchannel word $\rightarrow$
\texttt{user\_backchannel}. This is what lets the companion distinguish ``the user
barged in'' from ``the user just said \emph{aham}'' from ``the user paused''.

\paragraph{Silence semantics.} The runtime inserts a \texttt{<|no voice|>} marker
on ASR silence, so the data must define what silence means in every state:
non-final chunk $+$ silence $\rightarrow$ keep speaking; \emph{final} chunk $+$
silence $\rightarrow$ \texttt{finish} (turn complete---the case the original data
never shows, since \texttt{finish} always directly follows the final chunk);
partial final chunk $+$ silence $\rightarrow$ keep speaking; and minimal
single-utterance records $[\text{user chunk}, \texttt{<|no voice|>}]$ at weight 5,
without which the model never fires \texttt{finish} on a two-item context.

\paragraph{Promptable-controller records.} To make the controller
\emph{promptable}, a system instruction (a \emph{rule}) is rendered as the first
context item and the boundary after a trigger chunk is labelled
\texttt{system\_take\_floor} (the system interrupts the user immediately) or
\texttt{system\_handover} (control passes to the big model). Synthetic scenarios
use the real dictation shape---digit chunks with commas, the full alphabet as
letters---so the companion learns the concept (number vs.\ letter, trigger vs.\
unrelated words) rather than memorized strings; negatives prevent spurious fires.
About half of the regular records additionally carry a generic instruction so the
base controller behaviours survive instruction-conditioned training.

\subsection{Loss and inference}
The loss is the same two-term objective as Section~\ref{sec:distill}: masked
cross-entropy with per-token weights at item boundaries (retuned for the 9-way
decision---finish $=3$, interruption $=6$, backchannel $=3$, thinking $=2$,
speaking $=2$, take\_floor $=5$, handover $=3$, \texttt{no\_tag} $=0.5$), plus a
per-record multiplier so rare controller patterns dominate the gradient, and a
forward-KL to the frozen base on content positions with $\lambda=0.3$. Preserving
the companion's language ability is not cosmetic---the rule mechanism is an
instruction-following behaviour.

At inference the companion is queried per event (new ASR words, silence timer,
assistant sentence boundary); a tag fires only if its probability exceeds a
per-tag threshold, and \texttt{user\_interruption} is additionally suppressed when
the last user item is the orchestrator-inserted \texttt{<|no voice|>} marker, so
silence can never hard-stop the assistant. The context rendering must match
training exactly: user items are always closed, even for partial content, and only
an in-progress assistant item is rendered open.

\section{Evaluation}
\label{sec:eval}

\subsection{Protocol metric (teacher-forced)}
Tag accuracy on a held-out duplex split: at each supervised tag position the
argmax must equal the supervised special token. This is the metric the three
approaches are usually compared on. We also report the mean per-token forward-KL
between each student and the frozen base on the held-out content positions---a
direct measure of ``everything else stayed the same''.

\subsection{Capability metrics}
18 pt-BR function-calling prompts scored by native Qwen tool calls (greedy,
tool-name plus required/expected arguments, digit-normalized) and 12 general
instruction/QA prompts. Standard benchmarks: HellaSwag, ARC-Challenge and MMLU
(2{,}000-example subsets), GSM8K (5-shot, 300), TruthfulQA (mc2) and IFEval (40)
via lm-evaluation-harness 0.4.12, raw prompts without a chat template and
identical configs across models. The companion is additionally evaluated on its
own tag-stream held-out split with per-tag precision/recall/F1 and on hand-built
runtime probes.

\subsection{Free generation: Full-Duplex-Bench}
\label{sec:fdb_setup}
The metrics above are teacher-forced: at every supervised position the model is
handed the correct prefix and only its next-token choice is scored. That is the
right measure of whether the protocol was \emph{learned}, but it says nothing
about whether the model \emph{uses} it when generating on its own, which is what a
deployed duplex system needs. We therefore add a free-generation evaluation
modelled on Full-Duplex-Bench (FDB)~\cite{lin2025fdb}.

\paragraph{Dimensions and metrics.} FDB v1.0 scores four behaviours:
\emph{pause handling} and \emph{backchannel} (the model must stay silent or briefly
acknowledge: TOR $\downarrow$, backchannel frequency $\uparrow$, JSD against a
ground-truth backchannel-time distribution $\downarrow$), \emph{smooth
turn-taking} (take the floor as soon as the user's turn ends: TOR $\uparrow$,
latency $\downarrow$) and \emph{user interruption} (yield, then answer the
interrupting turn: TOR $\uparrow$, latency $\downarrow$, and an LLM-judged 0--5
relevance of the post-interruption reply). We use FDB's thresholds verbatim---a
model utterance is a formal turn at ${\ge}1$\,s or ${>}3$ words, a backchannel at
${\le}3$\,s and ${\le}3$ words---and its interruption rubric verbatim. Each subset
is generated at FDB v1.0's published size (216/137/119/55/200 scenarios, 727 per
model), in pt-BR and in English.

\paragraph{Text-domain adaptation.} Our cascade decides turn-taking in text, so
the harness replays the ASR stream as timed events rather than audio and converts
text to seconds with a calibrated speaking rate (3.25 words/s, measured against
the pocket-tts voice the live bridge uses and anchored to the rate implied by
FDB's own released transcripts). Latency is the model's real inference compute
between the user's turn ending and the model committing---the honest text-domain
analogue of acoustic response latency, and the quantity that separates a
tag-emitting LLM from a separate policy companion. Interruption relevance is
judged by a local \texttt{Qwen3.5-9B} GGUF under FDB's own rubric; the judge is
not GPT-4o, so that column ranks our models against each other rather than against
the published absolute. All runs are greedy and every model is served through the
same llama.cpp stack at the same quantisation.

\paragraph{The input stream must match training.} The models never see whole
sentences. The micro-turn builder delivers each user turn as 1--7-\textbf{token}
pieces (Section~\ref{sec:shared}), so we re-use that algorithm verbatim, including
its mid-word cuts and leading-space artifacts, to generate the evaluated stream.
This is not a detail, and the ablation in Section~\ref{sec:ablation} separates the
two ways it matters. Our first harness deviated on two axes at once: it fed whole
sentences as ASR pieces \emph{and} joined several sentences into one turn. The
single-factor arms show that the second axis is what distorted the turn-taking
numbers---multi-sentence turns more than triple pause-handling TOR, because
sentence-final ambiguity scales with pieces per turn---while the first axis breaks
interruption recovery instead, because a one-piece interrupting turn is a shape the
training builder never produces. We therefore generate every task's turns and pieces
to the measured distribution of Section~\ref{sec:shared}.

\paragraph{Every scenario is scored.} FDB v1.0 consults its judge in exactly one
dimension, so the paper-comparable column is the 200-sample interruption subset---
that is the whole subset, not a subsample. Our harness additionally scores every
scenario in which the model spoke ($\sim$400 per run) on a shared 0--5 relevance
rubric, and scores the rest by rule (silence correct $\rightarrow$ 5, silence when
a reply was due $\rightarrow$ 0), giving a whole-dataset quality column.

\paragraph{The reference model.} To anchor the comparison we download the upstream
English \textsc{DuplexCascade} release, merge its PEFT state dict, quantize to
4-bit and serve it through the identical stack. Three conventions of its own
codebase had to be matched, and each is documented because getting them wrong
makes the reference look broken: it prepends BOS before the first micro-turn; it
\emph{speaks by emitting text} rather than a tag (its server streams any generated
text to TTS and uses the tags only to stop); and it commits to the floor on the
\emph{silence} that follows speech rather than on the final word, so scenarios
must include the trailing silence that real audio always has. Its tag vocabulary
also differs (``talking'' vs.\ ``speaking'', and no system-backchannel token), and
is carried per-model.

\section{Results}
\label{sec:results}

\subsection{Protocol vs.\ capability retention}
Table~\ref{tab:custom} is the central result for the two in-model approaches. The
distill model learns the protocol (96.4\% tag accuracy, 1.6\,pp behind SFT) while
keeping function calling and general instruction-following at \emph{base level}
(100\%), with a KL drift only slightly above the base floor. The plain-SFT
baseline learns the protocol slightly better (98\%) but regresses on both
format-following capabilities and drifts ${\sim}3.5\times$ further from the base
distribution.

\begin{table}[t]
\centering
\caption{Protocol and capability retention (our pt-BR eval sets).}
\label{tab:custom}
\begin{tabular}{lccc}
\toprule
 & base & \textbf{distill} & sft \\
\midrule
Tag accuracy (\%) & 0.0 & \textbf{96.4} & 98.0 \\
Function calling (\%) & 100 & \textbf{100} & 55.6 \\
General instruction/QA (\%) & 100 & \textbf{100} & 41.7 \\
KL drift vs base (nats) & 0.51 & 1.07 & 4.76 \\
\bottomrule
\end{tabular}
\end{table}

\subsection{Standard benchmarks}
Table~\ref{tab:bench} reports the lm-evaluation-harness results. The distill model
is within ${\sim}1$--2\,pp of base on knowledge, commonsense and reasoning,
\emph{above} base on factuality and instruction following, and at base on the
custom sets. The SFT baseline drops on MMLU, TruthfulQA, IFEval and the custom
sets. The one regression of the distill model is GSM8K, discussed in
Section~\ref{sec:discussion}.

\begin{table}[t]
\centering
\caption{Standard benchmarks (lm-evaluation-harness 0.4.12; subsets as noted).}
\label{tab:bench}
\footnotesize
\resizebox{\columnwidth}{!}{%
\begin{tabular}{lccc}
\toprule
Benchmark & base & \textbf{distill} & sft \\
\midrule
HellaSwag acc-norm ($n{=}2000$) & 0.584 & 0.566 & 0.593 \\
ARC-Challenge acc-norm ($n{=}2000$) & 0.584 & 0.563 & 0.574 \\
MMLU ($n{=}2000$) & 0.705 & 0.698 & \textbf{0.660} \\
GSM8K exact ($n{=}300$, 5-shot) & 0.890 & 0.777 & 0.827 \\
TruthfulQA mc2 & 0.626 & \textbf{0.649} & 0.601 \\
IFEval loose / strict ($n{=}40$) & 0.730/0.651 & \textbf{0.762/0.683} & 0.603/0.524 \\
\bottomrule
\end{tabular}%
}
\end{table}

\subsection{Companion: tag-stream accuracy and controller probes}
Table~\ref{tab:tags} reports the companion's per-tag results on its held-out split
(3{,}222 boundary decisions, restricted argmax over the tag vocabulary). The two
rare but critical controller decisions---barge-in interruption and backchannel
recognition---stay near-perfect, and the two promptable-controller tags are
perfect. The small accuracy drop versus an earlier revision is the price of the
instruction-conditioned training and the distill term.

\begin{table}[t]
\centering
\caption{Per-tag results of the trained companion on the held-out duplex split.}
\label{tab:tags}
\footnotesize
\begin{tabular}{lccc}
\toprule
\textbf{tag} & \textbf{P} & \textbf{R} & \textbf{F1} \\
\midrule
no tag (emit nothing) & 1.00 & 0.95 & 0.973 \\
user is speaking & 0.99 & 0.92 & 0.953 \\
user finish speaking & 0.79 & 0.98 & 0.877 \\
user interruption & 0.95 & 0.99 & 0.969 \\
user backchannel & 1.00 & 0.98 & 0.99 \\
user is thinking & 0.90 & 1.00 & 0.949 \\
system take floor & 1.00 & 1.00 & 1.00 \\
system handover & 1.00 & 1.00 & 1.00 \\
\midrule
accuracy & \multicolumn{3}{c}{0.956} \\
\bottomrule
\end{tabular}
\end{table}

Runtime probes on hand-built conversations exercise the query path with guards and
thresholds applied: a mid-question user chunk $\rightarrow$ keep waiting; a
finished question followed by silence $\rightarrow$ \texttt{finish} (fires the
answer); a barge-in during assistant speech $\rightarrow$ \texttt{interruption};
``aham'' during assistant speech $\rightarrow$ \texttt{backchannel} rather than an
interruption; silence during assistant speech $\rightarrow$ keep talking. All are
correct with probability ${\ge}0.83$ at the decision point, the remainder being
carried by the no-voice guard and the state machine. With a rule configured, a
letter inside a digit dictation fires \texttt{system\_take\_floor} at 0.60--0.86,
the same rule with digit chunks only does not fire, and without a rule the system
tags never fire at all.

\subsection{Free-generation turn-taking}
Table~\ref{tab:fdb} reports the FDB results for the three approaches, the untrained
base as a control, and the upstream release measured with the same harness. Table
\ref{tab:fdb2} gives the same comparison in English, which is out-of-distribution
for all three trained variants (\S\ref{sec:results_transfer}).

\begin{table*}[t]
\centering
\caption{Full-Duplex-Bench, pt-BR (727 scenarios per model). TOR and latency are
better when lower; TOR, relevance and quality are better when higher. ``Original''
is the upstream English release measured with the same harness; the published row
is its paper's English audio result and is a range check, not a baseline. Quality
is a whole-dataset 0--5 reply rubric (see \S\ref{sec:fdb_setup}).}
\label{tab:fdb}
\footnotesize
\resizebox{\textwidth}{!}{%
\begin{tabular}{lccccccccc}
\toprule
Model & Pause TOR\textsubscript{syn} & Pause TOR\textsubscript{cand} & BC TOR & Smooth TOR & Floor lat.\,(s) & Intr.\ TOR & Intr.\ relevance & Stop lat.\,(s) & \textbf{ATA} \\
\midrule
Companion (MTP-style)             & 0.073 & \textbf{0.000} & 1.000 & 0.975 & 0.33 & 0.990 & \textbf{4.36} & 0.066 & \textbf{0.973} \\
SFT (original recipe)             & 0.080 & 0.032 & 1.000 & \textbf{1.000} & \textbf{0.25} & \textbf{0.995} & 3.56 & \textbf{0.035} & \textbf{0.971} \\
Self-distillation                 & 0.095 & 0.065 & 1.000 & 0.983 & 3.29 & 0.995 & 1.59 & 0.351 & 0.955 \\
Original DuplexCascade (upstream) & 0.102 & 0.273 & \textbf{0.436} & 0.975 & 1.02 & 0.985 & 3.92 & 0.323 & 0.896 \\
Base (no fine-tune, control)      & 0.000 & 0.000 & 0.000 & 0.000 & -- & 0.000 & -- & 0.15 & 0.000 \\
\midrule
\emph{published DuplexCascade (EN audio)} & \emph{0.058} & \emph{0.222} & \emph{0.218} & \emph{0.832} & \emph{1.72} & \emph{0.955} & \emph{4.02} & \emph{1.23} & \emph{0.858} \\
\bottomrule
\end{tabular}%
}
\end{table*}

\begin{table*}[t]
\centering
\caption{Full-Duplex-Bench in English (same harness, same scenario sizes). English
is out-of-distribution for all four trained models---they are fine-tuned on pt-BR
only---so this measures transfer of the protocol, not language skill.}
\label{tab:fdb2}
\footnotesize
\resizebox{\textwidth}{!}{%
\begin{tabular}{lccccccccc}
\toprule
Model & Pause TOR\textsubscript{syn} & Pause TOR\textsubscript{cand} & BC TOR & Smooth TOR & Floor lat.\,(s) & Intr.\ TOR & Intr.\ relevance & Stop lat.\,(s) & \textbf{ATA} \\
\midrule
SFT (original recipe)             & \textbf{0.051} & \textbf{0.000} & 1.000 & \textbf{1.000} & \textbf{0.27} & 0.995 & 4.08 & \textbf{0.031} & \textbf{0.986} \\
Companion (MTP-style)             & 0.066 & \textbf{0.000} & 1.000 & \textbf{1.000} & 0.32 & \textbf{1.000} & \textbf{4.88} & 0.057 & \textbf{0.984} \\
Original DuplexCascade (upstream) & 0.095 & 0.102 & \textbf{0.345} & \textbf{1.000} & 1.16 & \textbf{1.000} & 4.35 & 0.251 & 0.951 \\
Self-distillation                 & 0.095 & 0.065 & 1.000 & 0.941 & 4.34 & 0.915 & 2.12 & 0.189 & 0.924 \\
Base (no fine-tune, control)      & 0.000 & 0.000 & 0.000 & 0.000 & -- & 0.000 & -- & 0.15 & 0.000 \\
\bottomrule
\end{tabular}%
}
\end{table*}

\paragraph{The harness reproduces the reference.} The upstream release, measured
here, lands on its own published numbers: pause handling 0.102/0.273 against the
published 0.058/0.222, smooth turn-taking 0.975 against 0.832, interruption TOR
0.985 against 0.955, interruption relevance 3.92 against 4.02, and ATA 0.896
against 0.858. A re-implementation that reproduces the reference model's scores is
measuring the same thing the reference measures, which is the strongest available
check on everything else in these tables.

\paragraph{All three approaches beat the reference on turn-taking accuracy.} They
reach ATA 0.973 (companion), 0.971 (SFT) and 0.955 (self-distillation)---every one
above the upstream model's 0.896 and above the published 0.858---and all three sit
at or better than the published figures on every TOR column. They get there
differently, and that difference is the practical point of the paper: the
companion has the best interruption relevance and the best candour-pause behaviour
while a separate controller is queried before audio can be cut (stop latency
0.066\,s), whereas SFT is an order of magnitude faster to stop and to respond
(0.035\,s and 0.28\,s) because one model does everything.

\paragraph{What teacher forcing hid.} The distill model is competitive on
turn-taking accuracy but has by far the weakest interruption relevance (1.59
against 3.56 and 4.36) and is $13\times$ slower to commit to the floor (3.29\,s
against 0.25\,s). Teacher-forced tag accuracy (96.4\%, Table~\ref{tab:custom},
1.6\,pp behind SFT) is blind to exactly this: both models put the right tag at the
right position when given the correct prefix, but in free generation the distill
model's continuation artifact---the known failure in which it opens by finishing
the user's own sentence---degrades the reply and delays the turn. The live bridge
masks it by stripping leading continuation punctuation; this evaluation does not,
so the cost is visible. This is the concrete form of the exposure-bias gap
between the protocol metric and deployed behaviour, and it is the strongest
argument for evaluating turn-taking in free generation rather than at supervised
positions.

\paragraph{Cross-lingual transfer.}
\label{sec:results_transfer}
All variants are fine-tuned on pt-BR only, and the protocol transfers: ATA is
0.986 (SFT), 0.984 (companion) and 0.924 (distill) on the English scenarios---SFT
and the companion are as good or better in English than in pt-BR, with the
companion's interruption relevance rising to 4.88. Only the distill variant
degrades (0.955$\rightarrow$0.924, interruption TOR 0.995$\rightarrow$0.915),
consistent with it being the most fragile of the three. The ranking of the
approaches is stable across languages, which is a good sign that it reflects the
models rather than the run.

\paragraph{Backchannel is a data gap, not a capability gap.} Every pt-BR variant
scores BC TOR $=1.000$ where the upstream model scores 0.436: they always take a
formal turn during a monologue, while upstream backchannels. The cause is in the
training data---across 10{,}249 sequences the tag
\texttt{<|system backchannel|>} occurs \textbf{zero} times, and
\texttt{<|user backchannel|>} only 1{,}262 times (0.09\% of duplex tags)---so the
token exists in the vocabulary but was never trained. That column is not evidence
about these models and we do not interpret it. The same table also shows why a
single headline number is not enough: the untrained control scores
\texttt{chunk\_accuracy} 0.85 on the turn-end decision---\emph{higher} than every
trained model---purely by never taking the floor. Our averaged turn-taking
accuracy gates on the model demonstrably taking turns elsewhere and scores an
unmeasured interruption as zero, which puts the control at 0.000.

\subsection{Where the decisions go}
A diagnostic that turned out to explain the whole comparison: the two codebases
commit to the floor at \emph{different moments}. SFT, distill and the companion
fire on the final user chunk (rate 0.94--1.00) and essentially never on the
post-turn silence; the upstream release fires on the silence (0.95--0.99) and
almost never on the final chunk (0.013). Both conventions reach ${\sim}1.0$ smooth
turn-taking and interruption TOR---they are inverses of one another, not better or
worse---and this is why the upstream model measured as unable to take a turn at
all until its trailing-silence cue was supplied. Meanwhile
\texttt{pause\_stayed\_silent} is ${\sim}1.000$ for every trained model, original
included: \emph{nobody} speaks during a mid-turn pause. The differences between
these models are about \emph{when} they commit, not about whether they can hold.

\section{Ablations}
\label{sec:ablation}

We isolate the three factors that the results above depend on. Each arm changes
exactly one factor with respect to the main run---same 727 scenarios, seed, clock,
decoding and serving stack---so the differences in Table~\ref{tab:abl} are
attributable to that factor alone.

\begin{table*}[t]
\centering
\caption{Ablations on the pt-BR FDB harness. Each arm differs from its main run in
one factor. A1/A2 ablate the SFT model; A3 ablates the companion. TOR $\downarrow$
for the pause and backchannel columns, $\uparrow$ for smooth and interruption.}
\label{tab:abl}
\footnotesize
\resizebox{\textwidth}{!}{%
\begin{tabular}{llccccc}
\toprule
Arm & Factor changed & Pause TOR\textsubscript{syn} & Pause TOR\textsubscript{cand} & Smooth TOR & Intr.\ TOR & \textbf{ATA} \\
\midrule
\multicolumn{7}{l}{\emph{SFT: input-stream and turn-length axes}}\\
main            & none (training pieces, training-sized turns) & 0.080 & 0.032 & 1.000 & 0.995 & \textbf{0.971} \\
A1              & user pieces $=$ whole sentences & 0.080 & 0.097 & 1.000 & \textbf{0.000} & 0.706 \\
A2              & user turns $=$ multi-sentence (2--4) & \textbf{0.307} & 0.056 & 1.000 & 0.995 & 0.908 \\
\midrule
\multicolumn{7}{l}{\emph{Companion: controller augmentations}}\\
main (v6)       & full augmentation set, 9 tags & \textbf{0.073} & \textbf{0.000} & 0.975 & 0.990 & \textbf{0.973} \\
A3 (v5)         & no promptable-controller set, 7 tags & 0.088 & 0.097 & \textbf{1.000} & \textbf{1.000} & 0.954 \\
\bottomrule
\end{tabular}%
}
\end{table*}

\paragraph{Turn length, not piece granularity, drives pause handling.} The two
axes of the input distribution had to be separated, and doing so corrected our own
earlier account of the harness bug. Feeding whole sentences as ASR pieces (A1)
leaves pause handling essentially unchanged---0.080 synthetic, 0.097 candour---while
making the turns multi-sentence (A2) more than triples the synthetic-pause TOR to
\textbf{0.307}. The cause is the number of sentence-final pieces per turn: the real
data carries ${\sim}1.5$, A2 carries $3.1$, and each additional sentence-final piece
is another opportunity to fire \texttt{finish} early. What moves the pause metric
is therefore \emph{how much of the turn is still unheard}, not how the ASR slices
it. Our first harness did both at once and we initially attributed the effect to
the slicing; the single-factor arms show it belongs to the turn length.

\paragraph{Piece granularity drives interruption recovery.} The opposite holds for
interruptions: A1 collapses interruption TOR from 0.995 to \textbf{0.000} while
leaving every other column intact. A one-piece interrupting turn is
off-distribution in a specific way---the training builder refuses to interrupt with
a single-chunk turn, because that would produce two consecutive assistant messages
and break strict alternation---so the model has never seen an interruption delivered
as one long piece and does not re-take the floor afterwards. This is a concrete
example of a training-data invariant that a benchmark can silently violate, and it
is why the harness reproduces the builder's piece sizes rather than approximating
them.

\paragraph{Controller augmentations pay for themselves in precision.} The companion
ablates cleanly against its earlier revision (A3): without the promptable-controller
records and the two extra tags, ATA falls from 0.973 to 0.954. The trade is
instructive---A3 has \emph{better} recall on the turn-end decision (1.000 vs.\ 0.958)
but worse precision (0.680 vs.\ 0.698), and that precision loss is what shows up as
worse pause handling (0.088/0.097 vs.\ 0.073/0.000). The augmentations make the
controller more conservative about committing, which is exactly what the
stay-silent tasks reward. Almost all of the missing behaviour is in the candour-pause
column; the interruption and smooth-turn columns are unchanged or slightly better in
the ablation, so the cost of the extra machinery is confined rather than global.

\paragraph{The clock is not a sensitive parameter.} Because durations in this
harness come from a calibrated speaking rate rather than from audio, we recomputed
every TOR column of the four trained models at 2.5, 3.0, 3.25 and 4.0 words/s---a
$\pm25\%$ band around the calibrated value. Every value is identical to three
decimals at all four rates. The reason is that the decisions sit far from the
duration gates: a six-word reply is 1.5--2.4\,s at these rates, comfortably above
the 1\,s ``formal turn'' threshold, so the word-count branch of FDB's rule decides
rather than the duration branch. The clock is therefore a documented approximation
that does not drive any conclusion in this paper. The latency columns are measured
inference compute and are unaffected by it by construction.

\section{Discussion and Limitations}
\label{sec:discussion}

\emph{Why the content-KL works.} The tags live at supervised positions and are
predicted with the original high weights, so they are learned essentially as
before. Every other position is anchored to the frozen base, which is exactly the
behaviour we want to preserve. The mechanism is a single-stage, per-token analogue
of the ``re-invoke'' two-stage procedure~\cite{tml2025} and of SDFT's on-policy
anchoring~\cite{selfdistill2026}, but it requires neither rollouts nor a reward.

\emph{Free-generation cost of the anchoring term.} The content-KL preserves
capabilities but, at $\lambda{=}1.0$, it also anchors the model to the teacher's
content \emph{generation style}, and the teacher was occasionally autocompleting
the user turn when its auxiliary trigger asked for a short answer. The
distillation therefore inherits that continuation artifact, which shows up in
Table~\ref{tab:fdb} as the worst interruption relevance and the slowest floor
commitment of any model measured, even though its teacher-forced tag accuracy is
96.4\%. The fix we planned is at the data level, not the loss: regenerate the
teacher content with a corrected trigger (``responda de forma natural e completa
\ldots\ não se repita'') and retrain, so the artifact is absent from the anchored
distribution rather than stripped at inference. The FDB harness now provides the
measurement that tells us whether that retrain succeeded---relevance and stop
latency, not tag accuracy.

\emph{The GSM8K caveat.} The distill model drops 11\,pp below base on GSM8K, a
long-form multi-step reasoning task. This is the honest tradeoff of a strong
anchoring term ($\lambda{=}1.0$): content distributions are heavily pinned to the
base, and on 300 samples the 4B base's fragile chain-of-thought is sensitive to
any distributional shift. Two natural follow-ups are a lower $\lambda$ (our sweep
did not include a math benchmark) and the EMA-teacher or on-policy variants, which
track the student's progress and may preserve adaptability on reasoning while
still reducing forgetting.

\emph{Turn size is an experimental variable.} Two generations of source dialogues
exist in this project---short-only with a fixed 10-token assistant micro-turn, and
the short/long/mixed mixture with 10--48-token chunking used here---and the
assistant-side length distribution has a heavy right tail (median 5 words, p90 12,
max 37). Because FDB's turn-taking gates are durations, that distribution moves
the metric: a model trained on longer bursts crosses the formal-turn threshold
with fewer words. Benchmarking a duration-gated protocol therefore requires
reporting the turn-length distribution on both sides, and we do so.

\emph{What the companion buys and what it costs.} The companion makes turn-taking a
reusable, attachable module: a 0.8B sidecar trained in hours on one GPU, swappable
across big models with no retraining, and it supports instruction-conditioned
interrupts that no single-model fine-tune provides. The costs are real: one extra
model in the loop (stop latency 0.066\,s versus 0.035\,s for SFT), a learned
turn-end signal that depends on the ASR punctuating its output---unpunctuated
transcripts delay the \texttt{finish} decision---and a controller whose
instruction-conditioned behaviours must be explicitly mixed with generic
instructions so the base behaviours survive.

\emph{Limitations.} Our free-generation harness is text-domain: durations come from
a calibrated speaking rate rather than from audio, latency is inference compute
rather than acoustic end-to-end latency, and the judge is a local 9B model, so the
relevance column anchors our models against each other but not against the
published GPT-4o absolute. The published row is English audio and our rows are
pt-BR text; we use it as a range check rather than a baseline. All results are from
a single seed per configuration, and the source dialogues are synthetic. Finally,
one FDB dimension is unanswerable on this data by construction, as discussed above.

\section{Conclusion}
We compared three ways to give a conversational LLM the duplex-cascade micro-turn
protocol, on one pt-BR pipeline and with one harness.

Plain \textbf{SFT} learns the protocol best (98\% tag accuracy) but regresses on
function calling and general instruction-following; it remains the fastest model at
runtime and, in free generation, the best all-round turn-taker of the three.
\textbf{Self-distillation} splits the objective by token position---weighted CE on
the tags, forward-KL to the frozen base everywhere else---and retains base-level
capabilities and knowledge without RL or rollouts, at no extra cost over SFT with
an offline teacher cache. It is competitive on turn-taking accuracy but inherits
the teacher's generation style, which costs reply relevance and latency in free
generation---a cost teacher-forced tag accuracy (96.4\%) cannot see. The
\textbf{MTP-style companion} separates the two skills instead of balancing them: a
0.8B sidecar predicts only the control stream and drives an untouched big model,
reaching 95.6\% tag-stream accuracy with 0.97/0.99 F1 on barge-in and backchannel
decisions and supporting instruction-conditioned interrupts, at the price of one
extra model in the loop.

Evaluated in free generation with Full-Duplex-Bench, all three beat the published
turn-taking accuracy, and re-measuring the upstream DuplexCascade release with the
same harness reproduces its published numbers---which is what makes the comparison
trustworthy. Three lessons generalize beyond this system: a protocol metric should
be paired with a free-generation metric whenever the protocol is \emph{used} rather
than merely reproduced; a generative benchmark inherits the training input
distribution, so that distribution is an experimental variable that must be
reported (our single-factor ablations move pause-handling TOR by $4\times$ and
interruption TOR to zero on the same scenarios); and a metric can be unanswerable by
construction, as with the backchannel dimension here, where the required tag never
appears in the training data. Per-token KL-anchoring to a frozen base remains a practical, RL-free
anti-forgetting mechanism; the work here makes its cost measurable and localizes
it to the teacher's generation style rather than its knowledge.

\section*{Use of Generative AI}
This manuscript was prepared with the assistance of generative artificial
intelligence. All experimental results and the associated code were reviewed and
verified by the author.

\begin{thebibliography}{99}

\bibitem{duplexcascade}
J.~Yang, Y.~Fujita, and Y.~Sudo,
``DuplexCascade: Full-duplex speech-to-speech dialogue with VAD-free cascaded
ASR--LLM--TTS pipeline and micro-turn optimization,'' arXiv:2603.09180, 2026.
Demo: sbintuitions.github.io/DuplexCascadeDemo.

\bibitem{selfdistill2026}
I.~Shenfeld, M.~Damani, J.~H{\"u}b{\"o}tter, and P.~Agrawal,
``On-policy self-distillation fine-tuning (SDFT): self-distillation enables
continual learning,'' arXiv:2601.19897, 2026.

\bibitem{lin2025fdb}
G.-T. Lin, J.~Lian, T.~Li, Q.~Wang, G.~Anumanchipalli, A.~H.~Liu, and H.-y.~Lee,
``Full-Duplex-Bench: A benchmark to evaluate full-duplex spoken dialogue models on
turn-taking capabilities,'' arXiv:2503.04721, 2025.

\bibitem{ross2011}
S.~Ross, G.~Gordon, and D.~Bagnell,
``A reduction of imitation learning and structured prediction to no-regret online
learning,'' in \emph{Proc. AISTATS}, 2011.

\bibitem{kirkpatrick2017}
J.~Kirkpatrick \emph{et al.},
``Overcoming catastrophic forgetting in neural networks,'' \emph{PNAS},
vol.~114, no.~13, pp.~3521--3526, 2017.

\bibitem{lifhoiem2017}
Z.~Li and D.~Hoiem,
``Learning without forgetting,'' \emph{IEEE TPAMI}, vol.~40, no.~12,
pp.~2935--2947, 2018.

\bibitem{mecklenburg2024}
L.~Mecklenburg \emph{et al.},
``Injecting new knowledge into large language models via supervised fine-tuning,''
arXiv:2404.00213, 2024.

\bibitem{chu2025}
T.~Chu, Y.~Zhai, J.~Yang, S.~Tong, S.~Xie, D.~Schuurmans, Q.~V.~Le, S.~Levine, and
Y.~Ma,
``SFT memorizes, RL generalizes: a comparative study of foundation model
post-training,'' in \emph{Proc. ICML}, 2025, pp.~10818--10838.

\bibitem{airef}
R.~Agarwal, N.~Vieillard, Y.~Zhou, P.~Sta{\'n}czyk, S.~Ramos, M.~Geist, and
O.~Bachem,
``On-policy distillation of language models: learning from self-generated
mistakes,'' in \emph{Proc. ICLR}, 2024, arXiv:2306.13649.

\bibitem{rlsrazor2025}
I.~Shenfeld, J.~Pari, and P.~Agrawal,
``RL's razor: why online reinforcement learning forgets less,''
in \emph{Proc. ICLR}, 2026, arXiv:2509.04259.

\bibitem{tml2025}
K.~Lu and Thinking Machines Lab,
``On-policy distillation,'' Thinking Machines Lab: Connectionism,
DOI 10.64434/tml.20251026, 2025.

\bibitem{bai2022}
Y.~Bai \emph{et al.},
``Constitutional AI: harmlessness from AI feedback,'' arXiv:2212.08073, 2022.

\bibitem{snell2022}
C.~Snell, D.~Klein, and R.~Zhong,
``Learning by distilling context,'' arXiv:2209.15189, 2022.

\bibitem{lora2022}
E.~J.~Hu \emph{et al.},
``LoRA: low-rank adaptation of large language models,'' in \emph{Proc. ICLR}, 2022.

\bibitem{qlora2023}
T.~Dettmers, A.~Pagnoni, A.~Holtzman, and L.~Zettlemoyer,
``QLoRA: efficient finetuning of quantized LLMs,'' in \emph{Proc. NeurIPS}, 2023.

\bibitem{deepseek}
DeepSeek-AI,
``DeepSeek-V3 technical report,'' arXiv:2412.19437, 2024.

\bibitem{qwen35}
Qwen~Team,
``Qwen3.5: towards native multimodal agents,'' qwen.ai, Feb.~2026.

\bibitem{unsloth}
Unsloth~AI, ``Unsloth: 2x faster free finetuning,'' https://unsloth.ai, 2026.

\end{thebibliography}

\end{document}
